\section{The tower of truncated totality}

We write $\mathcal T_n$ for the categories total for $W_n$. All categories and all sets of coefficients below are small. Actions are covariant. We first establish the formal reductions, then give a uniform family separating every level of the tower.

\begin{claim}{A.1}{Proved}{Action inheritance and initial cones}\label{A:cone}
Let $n\geq2$. Every action component of a category in $\mathcal Q_n$ belongs to $\mathcal Q_n$. If $E\in\mathcal Q_n$, adjoining a new initial object gives $E^{\triangleleft}\in\mathcal T_n$, and $E$ is a nonempty full cosieve in $E^{\triangleleft}$.
\end{claim}
\begin{proof}
Let $L:E\to\Set$, let $D$ be one component of $\El_E L$, and let $K:D\to\Set$. The formula
\[
 \widetilde K(e)=\coprod_{\substack{l\in L(e)\\(e,l)\in D}}K(e,l)
\]
defines a covariant action on $E$. For $a:e\to e'$, its map sends $(l,k)$ to $(L(a)l,K(a)k)$. Thus $\El_DK\to\El_E\widetilde K$, sending $((e,l),k)$ to $(e,(l,k))$ and retaining the underlying arrow, is an isomorphism of categories. This verifies the exact inheritance step of TOTASPH C.5 and PRETOT C.2: the quantifier remains over all actions, with no finiteness restriction.

For an action $L':E^{\triangleleft}\to\Set$, put $L=L'|_E$. Each $x\in L'(0)$ defines the section
\[
 s_x:E\longrightarrow\El_E L,\qquad e\longmapsto(e,L'(0\to e)x).
\]
Its image is full, isomorphic to $E$, and a cosieve: every arrow from one of its objects has its target in the same section. Since $E$ is connected, the image lies in one component $D$. A component of $\El_{E^{\triangleleft}}L'$ is obtained from such a $D$ by adjoining one source for each $x$ whose section lies in $D$, with exactly one arrow to each object of $s_x(E)$. On nerves this attaches the cones on these section copies. Its realization is the homotopy pushout
\[
 BD\ \longleftarrow\ \coprod_x BE\ \longrightarrow\ \coprod_x *.
\]
Both $BE$ and $BD$ are $(n-1)$-connected. Van Kampen and the reduced Mayer--Vietoris sequence give simple connectivity and vanishing homology in degrees $1,\ldots,n-1$ for the new component, even for an infinite set of cones. Hurewicz gives $(n-1)$-connectedness. If there are no cones, the component is $D$ itself. The base $B(E^{\triangleleft})$ is contractible.

We use the category criterion of TOTASPH D.12 in exactly its stated form: $\mathcal T_n$ consists of the categories in $\mathcal Q_n$ whose base is $n$-connected. Its proof is valid: finality surjects on $\pi_1$; the comma homology spectral sequence is an isomorphism below $n$ and a surjection in degree $n$ when the commas are $(n-1)$-connected; conversely the collage forces homology vanishing in all action components, and the composite-action construction above permits the cyclic-cover argument to kill their fundamental groups. These are precisely the two directions used here. The criterion proves the assertion about the cone.
\end{proof}

\begin{claim}{A.2}{Proved}{The exact equality criterion behind PRETOT C.4}\label{A:equality-criterion}
For every $n\geq2$, the following conditions are equivalent:
\begin{enumerate}
\item Every nonempty full cosieve of a member of $\mathcal T_n$ belongs to $\mathcal T_n$.
\item $H_n(BE;\Z)=0$ for every $E\in\mathcal Q_n$.
\item $\mathcal Q_n=\mathcal Q_{n+1}$.
\item Totally $W_n$-dwy\'erienne sieves compose.
\end{enumerate}
Consequently these conditions hold for every $n\geq2$ if and only if every quasi-filtered category is total.
\end{claim}
\begin{proof}
For a nonempty cosieve $V\subset A\in\mathcal T_n$, extension of actions by the empty set off $V$ and PRETOT C.2 give $V\in\mathcal Q_n$. Its proof of connectedness uses the common-successor property forced by the vanishing of first homology of all actions; otherwise doubling $A$ along two disjoint upward closures produces a nonzero first-homology class. The criterion recalled in Claim~\ref{A:cone}, followed by Hurewicz, therefore gives
\[
 V\in\mathcal T_n\quad\Longleftrightarrow\quad H_n(BV;\Z)=0.
\]
Claim~\ref{A:cone} now proves (1)$\Leftrightarrow$(2). If (2) holds, apply it to each action component of $E\in\mathcal Q_n$, which is itself in $\mathcal Q_n$ by Claim~\ref{A:cone}. Each becomes $n$-connected, proving (3). Conversely apply (3) to the constant singleton action.

For completeness, we verify the exact step of TOTASPH F.9 used for (4). Given $U\subset X\subset Y$, at $y\in Y$ with $U/y\ne\varnothing$, the objects $(x\to y)$ with $U/x\ne\varnothing$ form a nonempty cosieve $S_y$ of $X/y$. If (1) holds, $S_y$ is total. The Grothendieck construction of $(x\to y)\mapsto U/x$ is total, by composition of totally aspherical functors, and its composition map to $U/y$ has the adjoint section $(u\to y)\mapsto(u\to y,u,\id_u)$. This map is final, so $U/y$ is total.

Conversely, for a nonempty cosieve $V\subset A\in\mathcal T_n$, form $X$ from an old copy of $A$ and a lower copy $V_0$, with
\[
 X(v_0,a)=A(v,a),\qquad X(a,v_0)=\varnothing,
\]
and all composites inherited from $A$. The sieve $V_0\to X$ is classical Dwyer: its generated cosieve is $V\times[1]$, whose projection to the lower copy is right adjoint to the inclusion. The copy-forgetting functor $X\to A$ is left adjoint to the old inclusion $A\to X$, so that inclusion is final and $X$ is total. Then $X\to X^{\triangleright}$ is totally dwy\'erienne; the boundary of the composite at the new terminal object is $V$. Composition would force $V$ total. This proves (4)$\Leftrightarrow$(1), with complete objects, arrows, and composition for the counterexample construction.

Finally $\mathcal Q_2$ is precisely quasi-filtered by Isbell--Mitchell's condition-I/affine-equalizer criterion, verified in R3, and $\bigcap_{n\geq2}\mathcal Q_n$ is the total class by the unrestricted-action criterion. Thus equality at every adjacent level is equivalent to their equality.

Here is the specific audit of R3 used in this last step and below. In TOTEX B.5, squares lift to a category of elements. An affine equalizer can also be lifted: partition its terms by their target element, equalize each part separately, and move the finitely many parts to one common successor. A loop then reduces by squares to a parallel pair, and the augmentation-one combination of its nerve triangles is an integral filling. Thus every action component has $H_1=0$. TOTEX B.6 applies this again to the cover associated with a nonzero cyclic subgroup; its $H_1$ would be that subgroup, proving simple connectivity. In the other direction the doubling-of-cosieves argument supplies common successors, applied also in each $p/A$ to supply squares, and the quotient representable for a parallel pair has $H_1=\Z/\operatorname{aug}\ker(P_q\to P_p)$. Its vanishing supplies one affine equalizer after moving finite supports to a common successor. This confirms both directions without invoking exactness of colimits for arbitrary abelian diagrams.
\end{proof}

\subsection{A uniform cocycle category}

\paragraph{Literature position.}
Campion--Ye realize arbitrary homotopy types by abstract elementary classes with amalgamation, in \cite[Theorem 3.6 and Corollary 3.7]{CampionYeAEC2020}. Their later cyclic-order calculation gives $K(\Z,2)$, in \cite[Corollaries 1.13--1.14 and 2.8]{CampionYeHigher2026}. The extra assertion needed here concerns every covariant set action; we prove that assertion explicitly below and in Claim~\ref{A:general-realization}. In degree two with $\mathbb F_2$ coefficients the finite structures are ordered two-graphs: the cocycle equation says that every four vertices span an even number of hyperedges, as in the definition of \cite{EvansHubickaKonecnyNesetril2020}. Their Corollary 7.2 shows that this ordered class is not Ramsey. Our proof uses the stated finite amalgams and open-subgroup calculation and assumes no Ramsey or imaginary-elimination theorem. No priority claim is made for these constructions or conclusions.

For a finite linearly ordered set $S$, let $\Delta^S$ be its full simplex. We use normalized simplicial cochains: an $r$-cochain with coefficients in an abelian group $\Lambda$ assigns a value to every increasing $(r+1)$-tuple; it extends alternatingly to distinct tuples and is zero on repeated entries. Its coboundary is the alternating sum omitting one entry. If $P\subset S$ and $p\in P$, the contraction
\[
 (s_pc)(v_0,\ldots,v_{r-1})=c(p,v_0,\ldots,v_{r-1})
\]
satisfies $d s_p+s_p d=\id$ in positive degrees.

\begin{claim}{A.3}{Proved}{Finite cocycle extension}\label{A:cocycle-extension}
Fix $r\geq2$ and an abelian group $\Lambda$. Let $K$ be a finite simplicial subcomplex of the full simplex on its vertices. If $H^r(K;\Lambda)=0$, every $r$-cocycle on $K$ extends to an $r$-cocycle on the full simplex. In particular this holds for a union of two simplices with a simplex as intersection (possibly empty), and for a union of three simplices whose nonempty intersections are simplices and whose nerve has dimension at most one or is a filled triangle.
\end{claim}
\begin{proof}
Write the given cocycle as $d b$ on $K$. Extend the $(r-1)$-cochain $b$ by zero on the other $(r-1)$-simplices; its coboundary is an extension of the cocycle. A union of two simplices as stated has no positive cohomology. For the three-piece assertion the elementary nerve lemma applies, since every nonempty intersection is contractible. The only case used below with all three pairwise intersections nonempty and empty triple intersection has nerve $S^1$. Its degree-$r$ cohomology is zero for $r\geq2$. Alternatively one obtains the same calculation by applying Mayer--Vietoris twice. This argument works over every abelian group; it is not a field-rank inference.
\end{proof}

\begin{claim}{A.4}{Proved}{Definition and homotopy type of the cocycle category}\label{A:cocycle-category}
For $r\geq2$ and a small abelian group $\Lambda$, define $\mathcal A_r(\Lambda)$ as follows. An object is a pair $(S,c)$, where $S$ is a finite nonempty linearly ordered set and $c\in Z^r(\Delta^S;\Lambda)$. An arrow $(S,c)\to(T,d)$ is an increasing injection $j:S\hookrightarrow T$ with $j^*d=c$. Composition is composition of injections. Taking $S=[m]$ gives a small skeleton. There is a specified weak equivalence
\[
 B\mathcal A_r(\Lambda)\simeq K(\Lambda,r).
\]
For $\Lambda\ne0$, its first nonzero homotopy and reduced integral homology groups are $\pi_r\cong H_r\cong\Lambda$.
\end{claim}
\begin{proof}
Let $K_m=Z^r(\Delta^{[m]};\Lambda)$. Pullback under all nondecreasing maps, with a cochain zero when a tuple repeats, makes $K$ a simplicial abelian group. Its normalized Moore group in degree $m$ is $\bigcap_{i=1}^m\ker d_i$. For $m<r$ it is zero; for $m=r$ it is $\Lambda$, the value on the full tuple. For $m=r+1$ a cochain in the intersection can be supported only on $\{1,\ldots,r+1\}$, and the cocycle equation on $\{0,\ldots,r+1\}$ forces that value to be zero. For $m>r+1$ no $(r+1)$-tuple contains every vertex $1,\ldots,m$, so the intersection is zero. Thus the normalized complex is $\Lambda[r]$. Moore's homotopy-group theorem for simplicial abelian groups identifies $K$ with an Eilenberg--Mac Lane space; see the opened Corollary III.2.7 of \cite{GoerssJardine1999}. Reversing the ordinal convention interchanges its normalization $\bigcap_{i=0}^{m-1}\ker d_i$ with the one used here. The Dold--Kan correspondence is Corollary III.2.3 of that source. We use only the homotopy-group theorem and the explicit normalized-complex calculation, not a claim about categories of nondegenerate simplices.

By definition $\mathcal A_r(\Lambda)=\Delta_{\mathrm{inj}}/K$, including every degenerate simplex as an object. The inclusion
\[
 j:\Delta_{\mathrm{inj}}/K\longrightarrow\Delta/K
\]
is a weak equivalence. Indeed, its overcategory at $\sigma:\Delta^m\to K$ consists of nondecreasing maps $[q]\to[m]$, with increasing injections between their domains that commute with these maps. Prepending an additional vertex labelled $0$ defines an endofunctor $T$ of this comma category, with natural arrows $\id\to T$ and from the constant one-vertex map labelled $0$ to $T$. It is therefore contractible. Quillen's Theorem A, in the overcategory variance, proves the assertion. Finally the standard simplex-category map $B(\Delta/K)\to|K|$ is a weak equivalence. It can be computed by the last-vertex map after simplicial replacement: the replacement has contractible simplices over each simplex of $K$. The Hurewicz theorem gives the claimed integral homology group.
\end{proof}

\begin{claim}{A.5}{Proved}{Functorial nonempty-base amalgamation}\label{A:cocycle-amalgam}
For $r\geq2$ and any abelian $\Lambda$, the category $\mathcal A_r(\Lambda)$ has joint embeddings and strong amalgamation. In every slice $P/\mathcal A_r(\Lambda)$ with $P\ne\varnothing$, every nonempty full cosieve has contractible classifying space.
\end{claim}
\begin{proof}
For two objects, or two objects sharing a subobject, choose a linear order on their set-theoretic amalgam extending the given orders. Their cocycles agree on their common simplex. Claim~\ref{A:cocycle-extension} extends them to the full simplex. The injections are induced embeddings and their images meet precisely in the prescribed base. This also supplies a joint embedding when the base is empty.

We require an actual functor for the second assertion. Fix $P\to Z$ in the cosieve and identify a variable object $P\to X$ with a tagged copy containing $P$. Put $p=\min P$. On $X$ and $Z$ take the canonical primitives $b_X=s_pc_X$ and $b_Z=s_pc_Z$. They restrict to the same primitive on $P$. On the tagged union $X\amalg_P Z$ define the $(r-1)$-cochain $b$ by $b_X$ on tuples wholly in $X$, by $b_Z$ on tuples wholly in $Z$, and by zero on tuples meeting both $X\setminus P$ and $Z\setminus P$. Give this union the cocycle $db$. In each gap of the order of $P$, place the $X$-side vertices before the $Z$-side vertices, preserving each side's order.

For an arrow $f:X\to X'$ under $P$, the map $f\amalg\id_Z$ is increasing and injective. It preserves the three cases in the definition of $b$, since injectivity and fixation of $P$ prevent a vertex outside $P$ from entering $P$. Pullback commutes with $s_p$ for a map fixing $p$, so it preserves $b$ and hence $db$. Thus $T(X)=X\amalg_P Z$ is a functor on the slice, with natural maps $\id\to T$ and $\mathrm{const}_Z\to T$. On the cosieve the image remains inside it, because it receives $Z\to T(X)$. These two natural transformations contract its nerve. Passing from tagged finite sets to the small ordinal skeleton transports this functor and its transformations along an equivalence; no functorial choice of an embedding into a limit is required.
\end{proof}

\begin{claim}{A.6}{Proved}{Finite supports for cocycle automorphisms}\label{A:cocycle-supports}
Let $r\geq2$ and let $\Lambda$ be countable. There is a countable ultrahomogeneous structure $M$ whose finite nonempty substructures are exactly the objects of $\mathcal A_r(\Lambda)$. Let $G=\operatorname{Aut}(M)$, with the pointwise convergence topology, and let $G_B$ denote the pointwise stabilizer of a finite subset $B\subset M$, including $G_\varnothing=G$. Then
\[
 \operatorname{Fix}_M(G_B)=B,\qquad
 \langle G_B,G_C\rangle=G_{B\cap C}.
\]
Every open subgroup of $G$ is $G_F$ for a finite set $F$, possibly empty.
\end{claim}
\begin{proof}
There are countably many finite cocycle structures and embeddings, because $\Lambda$ is countable. Heredity, joint embedding, and strong amalgamation were proved above. The usual explicit rich-chain construction therefore applies: enumerate all finite extensions of all earlier finite substructures, repeat each request, and amalgamate one requested extension at a time. The union $M$ has the extension property for every finite substructure. Back-and-forth on its countable underlying set extends every finite partial isomorphism to an automorphism. Its underlying order is countable dense without endpoints, since a new vertex can be inserted in any prescribed gap by strong amalgamation. If $x\notin B$, strongly amalgamating two copies of $B\cup\{x\}$ over $B$ and then using homogeneity supplies an automorphism fixing $B$ and moving $x$. This proves the fixed-point identity.

Put $D=B\cap C$. If one of $B,C$ is contained in the other (in particular, if either is empty), the generation identity is immediate from inclusion of stabilizers. We therefore assume neither is contained in the other. We prove the stabilizer-generation identity as an equality of subgroups, not merely after topological closure. Consider the bipartite incidence graph whose left vertices are copies of $B$ over $D$ under $G_D$, whose right vertices are copies of $C$ over $D$, and whose edges are copies of the whole ordered cocycle structure $B\cup C$ over $D$. By homogeneity this is the coset graph with vertices $G_D/G_B$ and $G_D/G_C$. Its connectedness is equivalent to $\langle G_B,G_C\rangle=G_D$: a finite path spells a finite alternating product from the two stabilizers, and conversely each such product traces a path.

We first treat two left vertices $B_0,B_1$ satisfying $B_0\cap B_1=D$. Suppose their corresponding order gaps overlap: for each gap in the ordered copy of $B$, the two open intervals determined by its images in $B_0$ and $B_1$ have a nonempty common interval whenever a vertex of $C\setminus D$ must be placed in that gap. We may then choose a new ordered copy $C'$ of $C$ over $D$, disjoint from $B_0\cup B_1$ off $D$, with the required order comparisons to both copies of $B$. Prescribe the cocycle on three simplices with vertex sets
\[
 B_0\cup B_1,\qquad B_0\cup C',\qquad B_1\cup C'.
\]
On the first use the existing cocycle of $M$; on the other two transport the cocycle of $B\cup C$. They agree on their overlaps $B_0,B_1,C'$, and their triple intersection is $D$. If $D$ is nonempty, the union has contractible nerve; if $D$ is empty, the nerve is a circle. Claim~\ref{A:cocycle-extension} applies because $r\geq2$. Extend to the full simplex and embed the finite extension into $M$ over $B_0\cup B_1$. The resulting $C'$ is adjacent to both $B_0$ and $B_1$.

For arbitrary copies $B_0,B_1$ over $D$, we insert a finite chain of intermediate copies, disjoint from each other and from $B_0\cup B_1$ away from $D$, such that neighboring corresponding gaps overlap. Here is a precise order construction. Identify the countable dense order of $M$ with a suborder of the real line and keep the coordinates of $D$ fixed. In each interval cut out by $D$, the space of strictly increasing tuples with the required number of coordinates is convex. Interpolate linearly between the tuples of $B_0\setminus D$ and $B_1\setminus D$. The finitely many coordinate gaps remain positive on this compact parameter interval. Choose a sufficiently fine finite mesh, and perturb its interior coordinates slightly so that they are all distinct and avoid the finitely many old coordinates. The successive corresponding gaps then overlap. This prescription is merely a finite total order pattern, so it can be realized in a finite extension.

To realize that order pattern together with the structures, prescribe the existing cocycle on the full simplex $B_0\cup B_1$ and the transported cocycle of $B$ on each intermediate copy. Distinct additional pieces meet only in $D$. Their union has no positive cohomology (it is a collection of simplices joined along the simplex $D$, or a disjoint union when $D$ is empty). The prescribed cocycle therefore extends to the full simplex. The extension property realizes all intermediate copies simultaneously in $M$ over the old union. Every neighboring pair is disjoint off $D$ and satisfies the overlap condition, so the preceding paragraph connects it by two edges. Thus all left vertices lie in one component; every right vertex has a left neighbor by homogeneity. The coset graph is connected, proving the identity.

Finally let $U\leq G$ be open. Some finite pointwise stabilizer lies in $U$; choose $F$ with $G_F\subset U$ and $|F|$ least. For $u\in U$, both $G_F$ and $G_{uF}$ lie in $U$, so the generation identity gives $G_{F\cap uF}\subset U$. Minimality forces $F\cap uF=F$, hence $uF=F$. A permutation preserving a finite linear order fixes every point, so $u\in G_F$. Consequently $U=G_F$, including the case $F=\varnothing$.
\end{proof}

\begin{claim}{A.7}{Proved}{Every action component is controlled}\label{A:all-actions}
Let $r\geq2$ and let $\Lambda$ be countable. For every action $L:\mathcal A_r(\Lambda)\to\Set$, each connected component $D$ of $\El L$ is either weakly contractible or its projection
\[
 D\longrightarrow\mathcal A_r(\Lambda)
\]
is a weak equivalence. In the second case $BD\simeq K(\Lambda,r)$. This assertion includes arbitrary infinite actions and arbitrary noninjective transition maps.
\end{claim}
\begin{proof}
We verify the all-action reduction used in TOTAL\_EXACT\_2 A.30--A.34, with its hypotheses supplied by Claims~\ref{A:cocycle-amalgam} and \ref{A:cocycle-supports}. Replace the ordinal skeleton by the equivalent concrete category $A$ of finite nonempty substructures of $M$, with all embeddings as arrows. This induces equivalences of categories of elements, so it changes neither assertion.

First remove noninjective transitions. On $L(a)$ define $l\sim l'$ when some $f:a\to b$ satisfies $L(f)l=L(f)l'$. Amalgamation of two witnessing arrows with source $a$ proves transitivity, and the same commuting-square argument proves preservation under every arrow out of $a$. Hence there is a quotient action $L^{\mathrm{inj}}$ whose transition maps are injective. Its element functor
\[
 q:\El_A L\longrightarrow\El_A L^{\mathrm{inj}}
\]
is homotopy final. For $(a,[l])$ in its target, define
\[
 J:a/A\longrightarrow (a,[l])\backslash q,
 \qquad f:a\to x\longmapsto ((x,L(f)l),f).
\]
For $d=((x,k),f)$ in the right-hand comma category, $d\backslash J$ is exactly the full cosieve in $x/A$ consisting of $g:x\to y$ such that $L(g)k=L(gf)l$. It is nonempty because $[k]=[L(f)l]$, and contractible by Claim~\ref{A:cocycle-amalgam}. Theorem A first makes $BJ$ a weak equivalence; since $a/A$ has an initial object, the target comma is contractible. A second application makes $Bq$ a weak equivalence. All these maps commute with projection to $A$.

For an injective action $L$, the finite subsets of $M$ with inclusion maps form a directed poset. Set $S=\colim_{B\subset M}L(B)$. Its canonical maps $i_B:L(B)\to S$ are injective. Define
\[
 g\,i_B(l)=i_{gB}(L(g|_B)l),\qquad g\in G.
\]
Compatibility with inclusions makes this a well-defined continuous discrete $G$-action: $G_B$ fixes every $i_B(l)$. For a continuous discrete $G$-set $S$, put $T_S(B)=S^{G_B}$. An embedding $f:B\to C$ extends to $g\in G$, and we set $T_S(f)s=gs$. Two extensions differ on the right by $G_B$, so this is independent of the choice; its value is fixed by $G_C$. Composition follows from extensions of composites. Thus $T_S$ is an action and $i:L\to T_S$ is a natural monomorphism.

For $(a,s)\in\El T_S$, the comma $(a,s)\backslash\El(i)$ is the full cosieve in $a/A$ of arrows $f:a\to b$ for which $T_S(f)s$ lies in $i_bL(b)$. The lifting label is unique. The cosieve is nonempty: represent $s$ at a finite $c$ and include $a\cup c$ in a finite $b$. Claim~\ref{A:cocycle-amalgam} makes it contractible, so $\El(i)$ is a weak equivalence, again over $A$.

Decompose $S$ into $G$-orbits. By Claim~\ref{A:cocycle-supports}, every orbit is $G/G_F$ for a finite $F$, possibly empty. If $F\ne\varnothing$, the natural identification
\[
 (G/G_F)^{G_B}\longrightarrow A(F,B),\qquad gG_F\longmapsto g|_F,
\]
is a bijection: a fixed coset means $G_B$ fixes $gF$ pointwise, hence $gF\subset B$ by the fixed-point identity; homogeneity proves surjectivity and equality of restrictions proves injectivity. It respects the transition maps. Its element category is $F/A$, which has an initial object. If $F=\varnothing$, $T_{G/G_F}$ is the constant singleton action and its element category is $A$ itself, with identity projection. Consequently each component of $\El T_S$ has exactly one of the two stated forms. The two weak equivalences over $A$ induce bijections on components and weak equivalences on every matched component. They transfer the conclusion back to the original arbitrary action $L$.
\end{proof}

\begin{claim}{A.8}{Proved}{Every adjacent level is different}\label{A:strict-tower}
For every $n\geq2$,
\[
 \mathcal A_n(\mathbb F_2)\in\mathcal Q_n\setminus\mathcal Q_{n+1},
 \qquad H_n(B\mathcal A_n(\mathbb F_2);\Z)\cong\mathbb F_2.
\]
In particular the separation asked in A.3 holds for every $n\geq3$.
There is a named nonzero integral cycle: let $t$ and $0$ be the two $n$-simplices of $K(\mathbb F_2,n)$, and let
\[
 \xi_n=\operatorname{Sd}_{\#}(t-0)\in C_n(N\mathcal A_n(\mathbb F_2);\Z).
\]
Here $\operatorname{Sd}_{\#}$ is the signed sum over flags of nonempty faces of the corresponding ordered simplex. Then $[\xi_n]\ne0$ and has order two.
\end{claim}
\begin{proof}
Claim~\ref{A:all-actions} controls every action component: it is contractible or a $K(\mathbb F_2,n)$, hence $(n-1)$-connected. Claim~\ref{A:cocycle-category} identifies the base as $K(\mathbb F_2,n)$. The base is itself the sole component of the constant singleton action, so the nonzero group in degree $n$ excludes $\mathcal Q_{n+1}$.

All proper faces of $t$ and $0$ agree, hence $t-0$ is a cycle in the unnormalized semisimplicial chain complex. Barycentric subdivision sends it to the displayed nerve cycle. More explicitly, for each permutation $\sigma$ of $0,\ldots,n$, use the flag
\[
 \{\sigma(0)\}\subset\{\sigma(0),\sigma(1)\}\subset\cdots\subset\{0,\ldots,n\}
\]
with coefficient $\operatorname{sgn}(\sigma)$, interpreted as the induced cocycle subobjects and increasing face injections; subtract the same flag for the zero cocycle. The normalized last-vertex cocycle with values in $\mathbb F_2$ evaluates to one on this chain: only the increasing permutation has strictly increasing successive maxima, and all other terms map to degenerate simplices. Thus $\xi_n$ is not an integral boundary. The Moore and Hurewicz calculations identify $H_n$ as $\mathbb F_2$, so its order is exactly two. The program accompanying this section checks the integer boundary and the mod-two evaluation independently for the first finite instances; the argument just given is uniform in $n$.
\end{proof}

\begin{claim}{A.9}{Refuted}{Cosieve closure and total dwy\'erienne composition at every finite level}\label{A:all-composition-failures}
For each $n\geq2$, nonempty full cosieves of categories total for $W_n$ need not be total for $W_n$, and totally $W_n$-dwy\'erienne sieves do not compose. The completely specified witness is
\[
 V=\mathcal A_n(\mathbb F_2),\qquad A=V^{\triangleleft}.
\]
Apply the lower-copy construction in Claim~\ref{A:equality-criterion} to obtain the two sieve inclusions and their failing composite. Moreover every inclusion in
\[
 \mathcal T_1\supsetneqq\mathcal Q_2\supsetneqq\mathcal T_2
 \supsetneqq\mathcal Q_3\supsetneqq\mathcal T_3
 \supsetneqq\mathcal Q_4\supsetneqq\cdots
\]
is strict, and the intersection is the total class $\mathcal Q_\infty$.
\end{claim}
\begin{proof}
Claims~\ref{A:cone} and \ref{A:strict-tower} give $A\in\mathcal T_n$ and $V\notin\mathcal T_n$, since $H_n(BV;\Z)\ne0$. The lower-copy construction then gives the two totally dwy\'erienne sieves with boundary $V$ for the composite. No finite search is being used to infer the all-action property.

The truncated criterion gives $\mathcal Q_{n+1}\subset\mathcal T_n\subset\mathcal Q_n$. The witness for $\mathcal Q_n\setminus\mathcal T_n$ is $\mathcal A_n(\mathbb F_2)$. The witness for $\mathcal T_n\setminus\mathcal Q_{n+1}$ is its initial cone: extend the singleton action on its cosieve by the empty set at the new initial object, so that its element category is exactly $V$. Finally $\mathcal T_1$ is the class with simply connected nonempty base, as verified by the finality-surjection argument in TOTASPH D.11. The three-object poset $0<1,0<2$ belongs to it and is not in $\mathcal Q_2$, because the two maximal objects have no common successor. All classes in the displayed chain contain the total categories; vanishing in every degree recovers the unrestricted-action criterion, proving the intersection assertion.

Thus the doubt at XII ms.~113--114 has a negative answer for totally $W_n$-dwy\'erienne sieves at every finite $n\geq2$, while the already established positive composition result for $W_\infty$ remains valid. This conclusion concerns sieve composition and does not change the positive composition theorem for strictly or totally aspherical functors.
\end{proof}

\subsection{The suspension criterion for arbitrary fundamental localizers}

\begin{claim}{A.10}{Proved}{Homotopy colimits for an arbitrary fundamental localizer}\label{A:localizer-hocolimits}
Let $W\supset W_\infty$ be a fundamental localizer, with no accessibility hypothesis. An objectwise $W$-equivalence of small category-valued diagrams induces a $W$-equivalence of their Grothendieck constructions. Consequently the induced class on spaces is preserved by the specified maps of homotopy colimits and, in particular, homotopy pushouts. We interpret a space $X$ by the category $\Delta/\operatorname{Sing}X$; this convention is independent of weak-equivalence replacements.
\end{claim}
\begin{proof}
For a covariant diagram $F:I\to\Cat$, let $p:\int_I F\to I$ be its opfibration. The category $(\int_I F)/i$ has a functor to $F(i)$ sending $(j,x,\alpha:j\to i)$ to $F(\alpha)x$. It is left adjoint to the section $y\mapsto(i,y,\id_i)$, hence a weak homotopy equivalence. An objectwise $W$-equivalence $F\to F'$ therefore gives $W$-equivalences on these comma categories, by two out of three, and the relative Theorem A axiom gives $\int F\to\int F'$ in $W$. This is the explicit special case of Cisinski's Proposition 3.3.16 used here; the statement is opened in \cite{Cisinski2006} and the proof above uses no model structure or small generating set. Opposite invariance of a fundamental localizer, as recorded in the same source, Sorites 3.3.8(7), gives the corresponding fibred statement.

For a simplicial-set diagram, apply the result to its simplex categories. The simplicial replacement has classifying space naturally weakly equivalent both to the nerve of the Grothendieck construction and to the homotopy colimit of the realizations. This comparison is obtained by contracting each individual simplex, so all its comparison arrows lie in $W_\infty$. It identifies the specified map induced by a natural transformation, not just the two abstract homotopy types. Taking the index to be the three-object span category proves the homotopy-pushout assertion. Applying singular complexes and their simplex categories to maps of spaces preserves identities, composition, and retract diagrams. Weak-equivalence replacements do not change membership in the induced class, by two out of three. These observations justify every homotopy-pushout comparison in the next proof without assuming accessible localization.
\end{proof}

\begin{claim}{A.11}{Proved}{Lead 6 in full generality}\label{A:suspension-criterion}
For every fundamental localizer $W\supset W_\infty$ and every nonempty small category $A$, the following are equivalent:
\begin{enumerate}
\item $A$ is total for $W$.
\item $BA$ is $W$-aspherical and, for every action $L:A\to\Set$ and every connected component $D$ of $\El_A L$, the unreduced suspension $\Sigma BD$ is $W$-aspherical.
\end{enumerate}
A categorical suspension model is $S(D)=D^{\triangleright}\amalg_D D^{\triangleright}$, with two distinct freely adjoined terminal objects and no arrows between them. Thus the statement refers to an explicit small category even when $W$ is not accessible.
\end{claim}
\begin{proof}
Suppose $A$ is total. Its identity final image makes $BA$ $W$-aspherical. For an action component $D$, restrict the action to its component subfunctor and form its collage $A_D$: there is one additional source $v$, and the arrows $v\to a$ are the elements of that subfunctor at $a$, with composition given by the action. The inclusion $A\to A_D$ is final. At an old object its outgoing comma has an initial object; at $v$ it is exactly $D$, which is nonempty and connected. Hence $A_D$ is $W$-aspherical. Its nerve is the mapping-cone realization of $BD\to BA$:
\[
 BA_D\simeq\operatorname{hocolim}(BA\longleftarrow BD\longrightarrow *).
\]
Indeed a new nerve simplex starts at $v$ and its remaining data are precisely a simplex of $ND$. By Claim~\ref{A:localizer-hocolimits}, the map obtained by replacing $BA$ with $*$ is a $W$-equivalence. Its target is $\Sigma BD$. Two out of three gives necessity.

Conversely let $f:A\to B$ be any final functor. For $b\in B$ put $D_b=b\backslash A$, contravariantly in $b$. Each $D_b$ is an action category of $A$ and is connected and nonempty, so $S(D_b)$ is $W$-aspherical by hypothesis. Form the fibred Grothendieck constructions over $B$ of $D$, of its two terminal cones, and of $S(D)$. Denote them by $E,E_+,E_-,E_S$, with their projections to $B$.

There is a functor $r:E\to A$ sending $(b,a,\beta:b\to f(a))$ to $a$. On an arrow it sends $(\delta:b\to b',\alpha:a\to a')$, satisfying $f(\alpha)\beta=\beta'\delta$, to $\alpha$. It is left adjoint to $s:A\to E$, $a\mapsto(f(a),a,\id)$: an arrow from $(b,a,\beta)$ to $s(a')$ is determined uniquely by $\alpha:a\to a'$, with base map $f(\alpha)\beta$. Thus $Br$ is a weak equivalence. The displayed $\beta$ form a natural transformation from the projection $p:E\to B$ to $fr$, so the map $BE\to BB$ corresponds to the specified map $Bf$ up to this stated homotopy. The two cone constructions have terminal objects in every fibre, preserved by transition functors. Their projections $E_\pm\to B$ have right adjoint sections and are weak equivalences.

The nerve of $E_S$ is the union of the two cone nerves along $NE$: after a chain reaches one of the new terminal sides, it remains on that side. The two inclusions are monomorphisms. Consequently the projection $E_S\to B$ identifies, through weak equivalences of the whole span, with the fold
\[
 d:BB\amalg_{BA}^{h}BB\longrightarrow BB.
\]
It lies in $W$, because its fibres $S(D_b)$ are $W$-aspherical and Claim~\ref{A:localizer-hocolimits}, or the fibred relative Theorem A calculation in its proof, applies.

Choose $a_0\in A$ and put $b_0=f(a_0)$. Since $BA\to *$ lies in $W$, its section $*\to BA$ at $a_0$ also lies in $W$. This gives an actual map of spans from
$BB\leftarrow *\to BB$ to $BB\leftarrow BA\to BB$, with identity maps on the two outside terms. Therefore the resulting map
\[
 BB\vee_{b_0}BB\longrightarrow BB\amalg_{BA}^{h}BB
\]
is in $W$. The composite with $d$ is the wedge fold, so this fold is in $W$. Each wedge inclusion $j_0,j_1:BB\to BB\vee BB$ is in $W$ by two out of three. Let $h:BB\vee BB\to BB$ kill the first summand at $b_0$ and fix the second. Since $hj_1=\id$, the map $h$ lies in $W$, and $hj_0$ is the constant self-map of $BB$ in $W$. Finally $BB\to *$ is a retract of that constant self-map in the arrow category, using $b_0:*\to BB$ and $BB\to *$ on the codomains. Retract closure gives $BB\to *\in W$. Since every final image $B$ is $W$-aspherical, $A$ is total. The only choice needed was $a_0$, available by the stated nonemptiness hypothesis.
\end{proof}

\begin{claim}{A.12}{Proved}{Suspensions, the tower, and strict intermediate conditions}\label{A:intermediate-localizers}
Write $\mathcal S_W$ for the nonempty connected categories whose every action component has $W$-aspherical suspension. Then
\[
 \mathcal S_{W_n}=\mathcal Q_n\quad(n\geq2),\qquad
 \mathcal S_{W_\infty}=\mathcal S_{W_{\Z}}=\mathcal Q_\infty.
\]
For $W_0$ and $W_1$, the suspension condition imposes no further condition on a nonempty connected category.

Fix $n\geq2$ and a fundamental localizer in the adjacent interval
$W_{n+1}\subseteq W\subseteq W_n$. Put
\[
 \mathcal R_W=\{\Lambda\in\Ab:K(\Lambda,n+1)\to *\text{ lies in }W\}.
\]
Then the exact condition is
\[
 A\in\mathcal S_W
 \quad\Longleftrightarrow\quad
 A\in\mathcal Q_n\ \text{ and }\
 H_n(BD;\Z)\in\mathcal R_W
 \text{ for every action component }D.
 \tag{A.12.1}
\]
In this interval the condition is strictly between $\mathcal Q_{n+1}$ and $\mathcal Q_n$ whenever $\mathcal R_W$ contains a nonzero countable group and omits a nonzero countable group. In particular strict intermediate conditions exist at every adjacent pair. One explicit choice is
\[
 W=W_n\cap W_{\Q,\leq n+1},
\]
where $W_{\Q,\leq n+1}$ denotes the functors inducing rational homology isomorphisms through degree $n+1$. For this choice (A.12.1) says that every $H_n(BD;\Z)$ is torsion. The categories $\mathcal A_n(\mathbb F_2)$ and $\mathcal A_n(\Z)$ witness the two strict inclusions.
\end{claim}
\begin{proof}
A suspension of a nonempty connected space is simply connected. Its reduced integral homology in degree $i+1$ is that of the original space in degree $i$. Thus the $W_n$ suspension condition says that every action component has zero reduced homology through degree $n-1$. Applied to cyclic-subgroup covering actions, available by the composite-action construction of Claim~\ref{A:cone}, its degree-one part forces every such component to be simply connected. Hurewicz then gives precisely $(n-1)$-connectedness. The same argument in every degree proves the two untruncated equalities. For $n=0,1$ the suspension of each connected component already has the required connectivity.

For the adjacent interval, membership in $\mathcal S_W$ first implies membership in $\mathcal Q_n$, since $W\subseteq W_n$. Conversely suppose $A\in\mathcal Q_n$. For each action component $D$, the space $\Sigma BD$ is $n$-connected, with $\pi_{n+1}=H_n(BD;\Z)$. Its Postnikov map
\[
 \Sigma BD\longrightarrow K(H_n(BD;\Z),n+1)
\]
is a $W_{n+1}$-equivalence, and hence is in $W$. Two out of three proves (A.12.1). It contains $\mathcal Q_{n+1}$ and lies in $\mathcal Q_n$. For a nonzero countable $\Lambda\in\mathcal R_W$, Claim~\ref{A:all-actions} shows $\mathcal A_n(\Lambda)\in\mathcal S_W\setminus\mathcal Q_{n+1}$. For a nonzero countable $\Lambda\notin\mathcal R_W$, it shows $\mathcal A_n(\Lambda)\in\mathcal Q_n\setminus\mathcal S_W$.

The class $W_{\Q,\leq n+1}$ is a fundamental localizer: weak saturation and the terminal-object axiom are immediate, and the same bisimplicial comma spectral sequence as for the homology localizer proves its relative Theorem A through the specified finite degree. Intersections of fundamental localizers are fundamental. Moreover $W_{n+1}\subseteq W_{\Q,\leq n+1}$, so this choice lies in the required interval. On an $n$-connected space the only possible reduced rational homology through degree $n+1$ is $H_{n+1}(-;\Z)\otimes\Q$. Hence $\mathcal R_W$ is exactly the torsion abelian groups. The displayed witnesses prove strictness.

For a general fundamental localizer outside an adjacent interval, Claim~\ref{A:suspension-criterion} still applies without change. Formula (A.12.1) expressly uses both interval inclusions; it is not asserted for an arbitrary localizer. Claim~\ref{A:global-intermediate} below gives the unrestricted classification in terms of ordinary covers and their suspensions, and removes the countability restriction from this interval's witnesses.
\end{proof}
